%% 
%% Copyright 2019-2021 Elsevier Ltd
%% 
%% This file is part of the 'CAS Bundle'.
%% --------------------------------------
%% 
%% It may be distributed under the conditions of the LaTeX Project Public
%% License, either version 1.2 of this license or (at your option) any
%% later version.  The latest version of this license is in
%%    http://www.latex-project.org/lppl.txt
%% and version 1.2 or later is part of all distributions of LaTeX
%% version 1999/12/01 or later.
%% 
%% The list of all files belonging to the 'CAS Bundle' is
%% given in the file `manifest.txt'.
%% 
%% Template article for cas-dc documentclass for 
%% double column output.
\documentclass[a4paper,fleqn]{cas-dc}
% If the frontmatter runs over more than one page
% use the longmktitle option.
%\documentclass[a4paper,fleqn,longmktitle]{cas-dc}

\usepackage[numbers]{natbib}
%\usepackage[authoryear]{natbib}
%\usepackage[authoryear,longnamesfirst]{natbib}
%\usepackage{amsmath}
%%%Author macros

\def\tsc#1{\csdef{#1}{\textsc{\lowercase{#1}}\xspace}}
\tsc{WGM}
\tsc{QE}
%%%
% Uncomment and use as if needed
\newtheorem{theorem}{Theorem}
\newtheorem{lemma}[theorem]{Lemma}
\newtheorem{proposition}{Proposition}
\newdefinition{rmk}{Remark}
\newproof{pf}{Proof}
%\newproof{pot}{Proof of Theorem \ref{thm}}
\newdefinition{definition}{Definition}
\newdefinition{example}{Example}

\let\WriteBookmarks\relax
\def\floatpagepagefraction{1}
\def\textpagefraction{.001}
% Short title
\shorttitle{Fractional Liu uncertain differential equation and its
	application to finance}    
% Short author
\shortauthors{Alireza Najafi, Rahman Taleghani}  
% Main title of the paper
\title [mode = title]{Comments for ‘‘Comments for Fractional Liu uncertain differential equation and its
	application to finance’’}  
% Title footnote mark
% eg: \tnotemark[1]
%\tnotemark[2] 
% Title footnote 1.
% eg: \tnotetext[1]{Title footnote text}
%\tnotetext[<tnote number>]{<tnote text>} 
% First author
%
% Options: Use if required
%\author[2]{Author Name}[type=editor,
%       style=chinese,
%       auid=000,
%       bioid=1,
%       prefix=Sir,
%       orcid=0000-0000-0000-0000
%      facebook=<facebook id>,
%       twitter=<twitter id>,
%       linkedin=<linkedin id>,
%       gplus=<gplus id>]

\author[1]{Alireza Najafi}

\cormark[*]
% Footnote of the first author
\fnmark[*]
% Email id of the first author
\ead{alireza.najafi@guilan.ac.ir}
% URL of the first author
%\ead[url]{https://staff.guilan.ac.ir/alireza.najafi/?lg=0}

% Credit authorship
% eg: \credit{Conceptualization of this study, Methodology, Software}
%\credit{<Credit authorship details>}

% Address/affiliation
\affiliation[1]{organization={Department of Applied Mathematics, Faculty of Mathematical Sciences, University of Guilan},
          %  addressline={}, 
            city={Rasht},
%          citysep={}, % Uncomment if no comma needed between     city and postcode
            postcode={41938-1914}, 
            state={Guilan},
            country={Iran}}

\author[1]{Rahman Taleghani}

% Footnote of the second author
%\fnmark[]

% Email id of the second author
\ead{Rahman.taleghani@gmail.com}

% URL of the second author
%\ead[url]{https://www.researchgate.net/profile/Rahman-Taleghani-2}
% Options: Use if required
%\author[2]{Rahman Taleghani}[type=editor,
%style=chinese,
%auid=000,
%bioid=1,
%prefix=[]
%orcid=0000-0002-5476-1927
%%      facebook=<facebook id>,
%%       twitter=<twitter id>,
%%       linkedin=<linkedin id>,
%%       gplus=<gplus id>
%]

% Credit authorship
%\credit{}
% Address/affiliation
%\affiliation[1]{organization={Department of Applied Mathematics, Faculty of Mathematical Sciences, University of Guilan},
	%addressline={}, 
%	city={Rasht},
	%          citysep={}, % Uncomment if no comma needed between city and postcode
%	postcode={41938-1914}, 
%	state={Guilan},
%	country={Iran}}
% Corresponding author text
\cortext[1]{Corresponding author}
% Footnote text
%\fntext[1]{}
% For a title note without a number/mark
%\nonumnote{}
% Here goes the abstract
\begin{document}


\begin{abstract}
In this article, we  demonstrate that the claim made by Xiaoyan Li and Ni Sun \cite{bib 1} regarding the incorrectness of Theorem 7 in the paper \cite{bib 2}  is wrong, and  show that this  Theorem is  based on the integral with respect to $(dt)^{\alpha}$  and the fractional derivative of order $\alpha$ ($ 0 < \alpha \leq 1$). 
\end{abstract}
% Use if graphical abstract is present
%\begin{graphicalabstract}
%\includegraphics{}
%\end{graphicalabstract}
% Research highlights
%\begin{highlights}
%\item 
%\item 
%\item 
%\end{highlights}
% Keywords
% Each keyword is seperated by \sep
\begin{keywords}
 Fractional derivative of order $\alpha$ \sep Integral with respect to $(dt)^{\alpha}$ \sep Fractional uncertain differential equation
 %%\MSC[2010] 60G22,  97M30.
\end{keywords}
\maketitle
% Main text


	
	
	%%author names are separated by comma (,) 
	%%use \and before the last author name 
	%%\textsuperscript{number} is used for affiliation
	%%use a * along with the number separated by comma
	%% for the  author for correspondence
	

	
	%%escape two column mode for title, affiliation and abstract
	%%by giving \twocolumn command as shown
	
	
	\section{Introduction}
	In 2005,  Guy Jumarie studied on the fractional Brownian motion and proposed a  representation of fractional Brownian motion as an integral with respect to $ (dt)^{\alpha}$. To do this, he used the fractional derivative of order $\alpha$ concepts \cite{bib 3}. After that, he  discussed about the  lagrange analytical mechanics and used the fractional derivative and the  integral with respect to $ (dt)^{\alpha}$ and   proposed an extension of the Lagrange analytical mechanics to deal with dynamics of fractal nature \cite{bib 4}. In 2010, Lv Longjin et al. studied about the  connection between fractional derivative and fractional Brownian motion  and used the  Jumarie's idea and presented an application of fractional derivatives in stochastic models driven by fractional Brownian motion \cite{bib 5}. In recent years, researchers used  the Jumarie's idea to find the option price PDE under the fractional  and mixed fractional financial models \cite{bib 6, bib 7, bib 8,bib 9}.
	
	
	
	In paper \cite{bib 2}, we used the integral with respect to $ (dt)^{\alpha}$  to show that the fractional uncertain differential equation has a unique solution.   Xiaoyan Li and Ni Sun in \cite{bib 1} by using the wrong definition of the  integral with respect to $ (dt)^{\alpha}$ proved that the Jumarie's idea and therefore  Theorem  7 that is  presented in \cite{bib 2} is wrong. Here, we study the Jumarie's idea and  a part of the   Theorem  7 that we used the idea. 
	
	
	
	
	
	\section{Main result}
	Here, we  present the integral with respect to $ (dt)^{\alpha}$ which is introduced by  Guy Jumarie and after that we mention some examples to better understand this definition.  In 2005,  Guy Jumarie created a bridge between the  integral with respect to $ (dt)^{\alpha}$ and the fractional integral as follows : 
	\begin{definition}\label{def 1}
		Let $f : \mathbb{R} \longrightarrow \mathbb{R}$ be a continuous function. Then its integral w.r.t $(dt)^{\alpha}$ for $0 < \alpha \leq 1$ is defined as follows:
		\begin{align}\label{eq 11}
			\int_0^t f(\tau) (d\tau)^{\alpha} = \alpha \int_0^t (t - \tau)^{\alpha -1}f(\tau)  d\tau
		\end{align}
	\end{definition}
	This definition is based on the  fractional derivative of order $\alpha$ for $0 < \alpha \leq 1$ \cite{bib 3}. Also, to better understand this definition, he gave some examples as follows \cite{bib 3, bib 4}:
	\begin{example}
		.
		\\
		(i) $~~$ Let $f(\tau) = 1$, then 
		\begin{align}
			\int_0^t  (d\tau)^{\alpha} = t^{\alpha} ~~~~ 0 < \alpha \leq 1.
		\end{align}
		(ii) $~~$ Let $f(\tau) =\tau^{\gamma}$, then 
		\begin{align}
			\int_0^t \tau^{\gamma} (d\tau)^{\alpha} = \frac{\Gamma(\alpha + 1) ~ \Gamma(\gamma + 1)}{\Gamma(\alpha + \gamma + 1)} t^{\alpha + \gamma}, ~~~~ 0 < \alpha \leq 1.
		\end{align}
		\\
		(iii) $~~$ Let $f(t)$ be  the Dirac delta generalized function $\delta (t)$, then 
		\begin{align}
			\int_0^t \delta (\tau) (d\tau)^{\alpha} = \alpha t^{\alpha -1}, ~~~~ 0 < \alpha \leq 1.
		\end{align}
	\end{example}
	It should be noted that this definition is based on fractional calculations and it cannot be expressed as a Riemann sum like the Riemann integral or Riemann Stieltjes Integral. 
	\\
	In paper \cite{bib 2}, we studied  the  fractional uncertain differential equation and proved that the equation has a unique solution. In fact, let $(\Gamma, \mathcal{M})$ be an uncertain space and  $\left\lbrace X_t \right\rbrace_{t \geq 0}$ satisfies in  
	\begin{align}\label{eq 43}
		dX_t = g(t, X_t) dt + h(t, X_t) dF_t^{\alpha},
	\end{align}  
	where $\left\lbrace F_t^{\alpha} \right\rbrace_{t \geq 0}$ is the fractional Liu process and   the coefficients $g(t, x)$ and $h(t, x)$ satisfy the linear growth and Lipschitz  conditions. Then the Eq \eqref{eq 43} has a unique solution. Here we present a part of the proof of theorem 7.  Assume that  $X_t^{(0)} = X_0$, and 
	\begin{align}
		X_t^{(n)} = X_0 &+ \int_0^t g\left( s, X_s^{(n-1)}\right) ds + \int_0^t h\left( s, X_s^{(n-1)}\right) dF_s^{\alpha},
	\end{align}
	for $n =1, 2, \cdots$, and suppose 
	\begin{align}
		D_t^{(n)}(\gamma) = \underset{0 < s < t}{max} \vert X_s^{(n+1)} (\gamma) - X_s^{(n)} (\gamma) \vert, 
	\end{align}
	for each $\gamma \in \Gamma$. It follows from the linear growth  and Lipschitz conditions and the fractional Liu process property, we have
	\begin{align}
	&	D_t^{(0)}(\gamma)  = \underset{0 < s < t}{max} \vert \int_0^s g\left(u, X_0\right) du + \int_0^s h\left(u, X_0\right) dF_u^{\alpha} \vert \nonumber \\
		& <  \int_0^t \vert g\left(u, X_0\right) \vert  du  +{\color{red} \frac{2 K(\gamma)}{1 - \alpha}\int_0^t \vert h\left(u, X_0\right) \vert  (du)^{1-\alpha}} \nonumber \\ 
		& = \int_0^t \vert g\left(u, X_0\right) \vert  du \nonumber \\ 
		&  + {\color{red}\frac{2 K(\gamma)}{1 - \alpha}(1-\alpha)\int_0^t (t - u)^{-\alpha} \vert h\left(u, X_0\right) \vert du} \nonumber \\
		& < L(1+|X_0|) \left(t + \frac{2 t^{1-\alpha} K(\gamma)}{1-\alpha}\right),
	\end{align}
	In this relation, we used Definition \ref{def 1} and  deduced that 
	\begin{align}\label{eq 26}
		\frac{2 K(\gamma)}{1 - \alpha}\int_0^t& \vert h\left(u, X_0\right) \vert   (du)^{1-\alpha} = \nonumber \\
		& \frac{2 K(\gamma)}{1 - \alpha}(1-\alpha)\int_0^t (t - u)^{-\alpha} \vert h\left(u, X_0\right) \vert du, 
	\end{align}
	where $h$ is a continuous function. \\
	In \cite{bib 1}, researchers used the Riemann sum and showed that the equation \eqref{eq 11} and thus the equation  \eqref{eq 26} are wrong. In fact, they used the Riemann sum concepts and considered the integral with respect to $(dt)^{\alpha}$ as follows 
	\begin{align}
		\int_0^t f(s) (ds)^{\alpha} = \underset{n \longrightarrow \infty}{\lim} \sum_{i=1}^n f(s_i) (\Delta s_i)^{\alpha},
	\end{align}
	where $0 = s_0 < s_1 < \cdots < s_n = t$. Then for $\alpha = \frac{1}{2}$, $t = 1$,  and $f(s) = 1$ for any $s \in [0, 1]$,  they concluded
	\begin{align}
		\int_0^1  (ds)^{\alpha} =  \underset{n \longrightarrow \infty}{\lim} \sum_{i=1}^n (\Delta s_i)^{\alpha}	= \underset{n \longrightarrow \infty}{\lim} \sum_{i=1}^n \left(\frac{1}{n}\right)^{\alpha} \longrightarrow  \infty
	\end{align}
	and $0 = s_0 < s_1 < \cdots < s_n = 1$. Therefore,  the $\int_0^1  (ds)^{\alpha}$ is diverge, and thus  the definition \ref{def 1}  is meaningless. 
	
	
	
	Based on this argument, they  concluded that Theorem 7 in \cite{bib 2} is wrong and discussed another equation and considered it as a replacement for this equation.
	
	
	
	
	
	
	But according to the definition \ref{def 1}, we have 
	\begin{align}
		\int_0^1  (ds)^{\alpha} = \alpha \int_0^1  (1 - s)^{\alpha -1} ds = 1. 
	\end{align}
	
	
	
	
	
	
	
	
	
	
	
	
	
	
	
	
	%%use \balance somewhere in the last page to balance the 2 columns.
	
	%%References section
	\begin{thebibliography}{99} 
		
		\bibitem{bib 1}
		Li, X., \& Sun, N. (2023). Comments for “Fractional Liu uncertain differential equation and its application to finance”. Chaos, Solitons \& Fractals, 175, 113293.
		\bibitem{bib 2}
		Najafi, A., \& Taleghani, R. (2022). Fractional Liu uncertain differential equation and its application to finance. Chaos, Solitons \& Fractals, 165, 112875.
		\bibitem{bib 3}
		Jumarie, G. (2005). On the solution of the stochastic differential equation of exponential growth driven by fractional Brownian motion. Applied Mathematics Letters, 18(7), 817-826.
		\bibitem{bib 4}
		Jumarie, G. (2007). Lagrangian mechanics of fractional order, Hamilton–Jacobi fractional PDE and Taylor’s series of nondifferentiable functions. Chaos, Solitons \& Fractals, 32(3), 969-987.
		\bibitem{bib 5}
		Longjin, L., Ren, F. Y., \& Qiu, W. Y. (2010). The application of fractional derivatives in stochastic models driven by fractional Brownian motion. Physica A: Statistical Mechanics and its Applications, 389(21), 4809-4818.
		\bibitem{bib 6}
		Jumarie, G. (2008). Stock exchange fractional dynamics defined as fractional exponential growth driven by (usual) Gaussian white noise. Application to fractional Black–Scholes equations. Insurance: Mathematics and Economics, 42(1), 271-287.
		\bibitem{bib 7}
		Guo, Z., \& Yuan, H. (2014). Pricing European option under the time-changed mixed Brownian-fractional Brownian model. Physica A: Statistical Mechanics and its Applications, 406, 73-79. 
		\bibitem{bib 8}
		Murwaningtyas, C. E., Kartiko, S. H., Gunardi, H. P. S., \& Suryawan, H. P. (2020). Finite Difference Method for Pricing of Indonesian Option under a Mixed Fractional Brownian Motion. Mathematics and Statistics, 8(5), 
		610-619.
		\bibitem{bib 9}
		Najafi, A., \& Mehrdoust, F. (2023). Conditional expectation strategy under the long memory Heston stochastic volatility model. Communications in Statistics-Simulation and Computation, 1-21.
	\end{thebibliography}
	
\end{document}

%
%% Biography
%\bio{}
%% Here goes the biography details.
%\endbio
%
%%\bio{pic1}
%% Here goes the biography details.
%\endbio


